\subsection{Visual Design}\label{subsect:VisualDesign}

The visual design follows the component sheet created in Figma. Inter is used as the primary typeface, with a clear distinction between page titles, section titles, body content, and secondary labels. Cards use generous rounded corners and light borders, while primary actions use a pill shape. Status information is never conveyed only by color: labels, icons, and progress values provide the same meaning in text.

The palette is stored as design tokens rather than repeated manually across screens. The main brand color is \texttt{\#FE6363}, the application background is \texttt{\#F9FAFB}, and additional semantic tokens represent success, warning, danger, text, borders, and appointment categories. The same token file used by the Figma project is imported by the application theme, reducing the possibility of the prototype and implementation drifting apart.

Material UI's \texttt{createTheme} function defines typography, spacing, shape, palette, and component overrides in one location. Responsive font sizing is then applied so headings adapt across mobile and wider layouts. Appointment cards use an additional semantic palette group, \texttt{palette.task}, with one entry per appointment kind and per terminal state: training, protocol, nutrition, done, and suspended. The first four come from the token file's task group and the last from its status group, so a card's colour is decided in Figma and read in the theme rather than chosen again in the component. All screen-level styling is expressed through Material UI and the theme rather than through separate CSS files.

The final interface preserves the soft cards and compact navigation of the prototype while improving clarity in nested pages. Shared page headers provide a consistent back action, sticky top and bottom regions keep global navigation available, and forms place validation or offline information next to the action it affects. Notification filters and selection states use labels and checkboxes in addition to color, while the membership banner keeps an account-wide restriction visible without hiding the page underneath it.
